% =============================================================================
% CAPÍTULO 1: Introducción al Trabajo Fin de Máster
% =============================================================================
\chapterUnnumbered{Capítulo 1: Introducción al Trabajo Fin de Máster}

[Todo lo que veas entre corchetes son explicaciones a borrar de la
versión final de tu TFM (Corchetes incluidos).

El TFM se distribuye en Capítulos > Apartados > Subapartados. Si
necesitas un cuarto nivel puedes añadir Subapartados de 2º nivel con el
mismo formato que los subapartados de 1\textsuperscript{er} nivel.

Redacta utilizando el estilo que corresponda de los 8 disponibles en la
plantilla (en la pestaña INICIO arriba en el encabezado de Word tienes
todos los estilos disponibles), así siempre seguirás respetando el
formato oficial de la UAX.

No utilices sangrías ni cambies el estilo de texto porque ya está
adaptado al formato UAX. Cuando pongas títulos usa el estilo \textbf{Título
Capítulo} (título de capítulo de primer nivel), \textbf{Título Apartado}
(para los apartados en los que se subdivide el capítulo) y \textbf{Título
Subapartado} (para los subapartados en los que se subdivide el
apartado) según proceda, \textbf{Texto TFM} para el texto, \textbf{Bibliografía},
para la bibliografía, y \textbf{Cita Literal APA} para las citas literales de
más de 40 palabras. Las viñetas son \textbf{Viñetas TFM} y las \textbf{Viñetas
numeradas TFM}. Si no sabes lo que es una viñeta fíjate en esto:]

\begin{itemize}
  \item Hola, yo soy una viñeta :)
\end{itemize}

\begin{enumerate}
  \item Hola, yo soy una viñeta numerada XD
\end{enumerate}

CONSEJO GENERAL DE REDACCIÓN: Recuerda que la forma más adecuada de
redactar es con estilo impersonal, es decir: Se ha realizado...; se
parte de...; Se pretende desarrollar... En vez de: Realicé/realizamos...
Partimos de... Pretendí/pretendimos realizar...

Antes de comenzar a redactar el primer apartado del primer capítulo
(CONSEJO GENERAL: Evita siempre poner 2 títulos seguidos sin estar
separado por texto) en un primer párrafo describe de forma simple en qué
consiste tu trabajo para que el lector sepa que se va a encontrar en el
TFM. En un segundo párrafo explica brevemente los contenidos de este
capítulo antes de continuar con los siguientes apartados (CONSEJO
GENERAL: Cada vez que de un capítulo se deriven varios apartados o de
que de un apartado se deriven varios subapartados conviene mostrar en el
último párrafo antes de la derivación cuáles son dichas partes en las
que se deriva el capítulo/apartado).

\setcounter{section}{0}
\section{Justificación del proyecto realizado}

[Primero justificaremos de manera personal por qué elegimos este
proyecto (¿Cómo se te ocurrió la realización de este proyecto? ¿Por qué
crees que lo vas a realizar bien?), y luego justificaremos la utilidad
científica/educativa del proyecto (¿Por qué crees que será útil? ¿Por
qué crees que es relevante realizarlo como TFM?).

Seguramente necesites contextualizar el tema de estudio en la actualidad
y en el contexto donde se va a aplicar de forma muy general. Es decir,
habla de lo que dicen los ámbitos científico, social y educativo hoy en
día respecto al tema tratado sin mucha profundidad (CONSEJO GENERAL:
Dedícale un pequeño párrafo a describir cada capítulo puede ser
efectivo).]

\section{Presentación del proyecto y sus contenidos}

[Describe de forma simple el problema que presentas y cómo tu proyecto
contribuye a resolverlo. Posteriormente, haz un recorrido por la
distribución de los contenidos en tu proyecto describiendo de forma
simple los capítulos y apartados que vas a incluir, de forma lógica,
para preparar al lector a su lectura.]

[Se recomienda un tamaño de 5 a 10 páginas para este capítulo].
